% Opcoes editoriais de trabalho, nao medidas impostas pela revista:
% A4, 12 pt, margens de 2,5 cm, entrelinha 1,5, Latin Modern.
\usepackage[T1]{fontenc}
\usepackage[utf8]{inputenc}
\usepackage{lmodern}
\IfFileExists{portuguese.ldf}{%
  \usepackage[portuguese,english]{babel}
}{%
  \usepackage[english]{babel}
  % Se os padroes portugueses nao estiverem instalados, nao aplicar
  % hifenizacao inglesa aos trechos em portugues.
  \babelprovide[import,hyphenrules=nohyphenation]{portuguese}
}
\usepackage[margin=2.5cm]{geometry}
\usepackage{microtype}
\usepackage{setspace}
\onehalfspacing
\usepackage[autostyle=false,style=english]{csquotes}
\usepackage{enumitem}
\usepackage{xcolor}
\usepackage{etoolbox}
\usepackage{xparse}
\usepackage{needspace}
\usepackage[most]{tcolorbox}
\usepackage{titlesec}
\usepackage{xurl}
\usepackage[hidelinks,unicode]{hyperref}
\urlstyle{same}
\hypersetup{
  pdftitle={How to talk about the Anthropocene? A Science-bar essay towards ancestral futures},
  pdfauthor={},
  pdfsubject={Essay manuscript},
  pdfkeywords={Anthropocene, Bruno Latour, common world, ancestral futures}
}

\setlength{\parindent}{1.25cm}
\setlength{\parskip}{0pt}
\setlength{\emergencystretch}{2em}
\clubpenalty=10000
\widowpenalty=10000
\displaywidowpenalty=10000
\raggedbottom
\titleformat{\section}{\large\bfseries}{\thesection.}{0.6em}{}
\titlespacing*{\section}{0pt}{2em}{0.8em}
\renewcommand{\thesection}{\Roman{section}}
\setlist[enumerate]{leftmargin=1.7em,itemsep=0.25em,topsep=0.4em}
\setlist[itemize]{leftmargin=1.4em,itemsep=0.25em,topsep=0.4em}

% Conteudo editorial: completamente removido no modo sem orientacoes.
\definecolor{guideline}{HTML}{BDC6CB}
\definecolor{guidefill}{HTML}{F4F6F7}
\NewDocumentEnvironment{orientacao}{m +b}{%
  \iforientacoes
    \begin{tcolorbox}[enhanced,breakable,colback=guidefill,colframe=guideline,
      boxrule=0.35pt,arc=0pt,left=9pt,right=9pt,top=7pt,bottom=7pt,
      before skip=1em,after skip=1em,
      title={#1},colbacktitle=guidefill,coltitle=black,
      fonttitle=\normalfont\sffamily\bfseries\small,
      fontupper=\normalfont\small]
      \selectlanguage{portuguese}\raggedright
      \setstretch{1.13}\setlength{\parindent}{0pt}\setlength{\parskip}{0.4em}
      #2
    \end{tcolorbox}
  \fi
}{}

% Atos: narrativa ficcional, sem aspas de citacao de fonte externa.
\newcounter{ato}[section]
\NewDocumentEnvironment{ato}{m}{%
  \refstepcounter{ato}
  \par\addvspace{1.2em}
  \noindent\begin{minipage}{\linewidth}
  \noindent\textbf{Act \Roman{ato}. #1}\par\smallskip
  \begingroup\itshape\setstretch{1.2}
  \setlength{\parindent}{0pt}\setlength{\parskip}{0.35em}
}{\par\endgroup\end{minipage}\par\medskip}

% Uma chave centraliza a nota curta e a referencia completa.
% O argumento opcional de \fonte recebe o localizador ja formatado.
\newcommand{\declararobra}[3]{%
  \csdef{curta@#1}{#2}\csdef{completa@#1}{#3}}
\newcommand{\usarobra}[1]{\csgdef{usada@#1}{1}}
\newcommand{\refobra}[1]{%
  \ifcsdef{curta@#1}{\usarobra{#1}\csuse{curta@#1}}{%
    \PackageError{sciencebar}{Referencia desconhecida: #1}{Cadastre a chave em referencias.tex.}}}
\NewDocumentCommand{\fonte}{O{} m}{%
  \footnote{\refobra{#2}\ifstrempty{#1}{}{, #1}.}}
\newcommand{\entradaobra}[1]{%
  \iforientacoes
    \par\noindent\hangindent=1.25cm\hangafter=1\csuse{completa@#1}\par\medskip
  \else
    \ifcsdef{usada@#1}{%
      \par\noindent\hangindent=1.25cm\hangafter=1\csuse{completa@#1}\par\medskip}{}
  \fi}

% As aspas sao visiveis e deliberadas; o texto explica sua funcao.
\newcommand{\anthropocene}{\enquote{Anthropocene}}
\newcommand{\sciencebar}{Science-bar}

% Espaco para atos ainda nao escritos nas partes II e III.
% Ao escrever, substituir \atoplanejado por \begin{ato} ... \end{ato}.
\newcommand{\atoplanejado}[2]{%
  \begin{orientacao}{Ato #1 --- a escrever}#2\end{orientacao}}
